\subsection{深度与零调性}

命题 3.--- 设 A 为环，C 为 A-模的下有界复形，p 为整数。假设对每对整数 \((m,n)\) （满足 \(m \geqslant n \geqslant p\) ），A-模 \(C_{m}\) 关于 \(\mathrm{H}_{n}(\mathrm{C})\) 的零化子的深度都 \textgreater m - n。则 \(\mathrm{H}_{n}(\mathrm{C}) = 0\) 对 \(n \geqslant p\) 成立。

由于 C 左侧有界，当 n 充分大时 \(\mathrm{H}_{n}(\mathrm{C})\) 为零。若结论不真，则存在整数 \(q \geqslant p\) ，使得 \(\mathrm{H}_{n}(\mathrm{C}) = 0\) 对 n \textgreater{} q 成立，而 \(\mathrm{H}_{q}(\mathrm{C}) \neq 0\) 。以 J 表示 \(\mathrm{H}_{q}(\mathrm{C})\) 的零化子；则有 \(\operatorname{prof}_{\mathrm{A}}(\mathrm{J}; \mathrm{H}_{q}(\mathrm{C})) = 0\) 。此外，由于 \(\mathrm{Z}_{q}(\mathrm{C})\) 是 \(C_{q}\) 的子模，且由假设 \(\operatorname{prof}_{\mathrm{A}}(\mathrm{J}; \mathrm{C}_{q}) > q - q = 0\) ，故有 \(\operatorname{prof}_{\mathrm{A}}(\mathrm{J}; \mathrm{Z}_{q}(\mathrm{C})) > 0\) 。由正合列

\[
0 \to \mathrm{B} _ {q} (\mathrm{C}) \to \mathrm{Z} _ {q} (\mathrm{C}) \to \mathrm{H} _ {q} (\mathrm{C}) \to 0
\]

导出等式 \(\mathrm{prof}_{\mathrm{A}}(\mathrm{J};\mathrm{B}_q(\mathrm{C})) = 1\) （第 1 节, 命题 1）。按 \(q\) 的定义，\(\mathrm{B}_n(\mathrm{C})\) 等于 \(\mathrm{Z}_n(\mathrm{C})\) ，对每个整数 \(n > q\) 皆然。由典范正合列

\[
0 \to \mathrm{B} _ {n} (\mathrm{C}) \to \mathrm{C} _ {n} \to \mathrm{B} _ {n - 1} (\mathrm{C}) \to 0 \quad (n > q)
\]

以及假设 \(\operatorname{prof}_{\mathrm{A}}(\mathrm{J};\mathrm{C}_{n}) > n - q\) ，用归纳法得到等式 \(\operatorname{prof}_{\mathrm{A}}(\mathrm{J};\mathrm{B}_{n}(\mathrm{C})) = n - q + 1\) 对每个 \(n \geqslant q\) 成立（同上引文）。但这是荒谬的，因为 \(\mathrm{B}_{n}(\mathrm{C})\) 当 \(n\) 充分大时为零。

推论 1.--- 设 A 为环，J 为 A 的理想，C 为 A-模的下有界复形，p 为整数。假设 \(\mathrm{JH}_{m}(\mathrm{C}) = 0\) 且 \(\operatorname{prof}_{\mathrm{A}}(\mathrm{J}; \mathrm{C}_{m}) > m - p\) 对 \(m \geqslant p\) 成立。则 \(\mathrm{H}_{n}(\mathrm{C}) = 0\) 对 \(n \geqslant p\) 成立。

事实上，对 \(n \geqslant p\) ，零化子 \(J_n\) （即 \(H_n(C)\) 的零化子）包含 \(J\) ；因此有 \(\mathrm{prof}_{\mathrm{A}}(\mathrm{J}_{n};\mathrm{C}_{m}) \geqslant \mathrm{prof}_{\mathrm{A}}(\mathrm{J};\mathrm{C}_{m})\) （ \(n^{\circ}\) 1, 命题 2 的推论 1），从而命题的假设得到满足。

推论 2.-- 设 A 为局部环，C 为 A-模的下有界复形，p 为整数。假设对 \(m \geqslant p\) ，\(\mathrm{H}_{m}(\mathrm{C})\) 的长度有限且 \(\mathrm{C}_{m}\) 的深度 \(>m-p\) 。则 \(\mathrm{H}_{n}(\mathrm{C}) = 0\) 对 \(n \geqslant p\) 成立。

A-模 \(\bigoplus_{m\geqslant p}\mathrm{H}_{m}(\mathrm{C})\) 的长度有限。以 J 表示其零化子；由 A, VIII, § 1, 第 3 节, 推论，环 A/J 是 Artin 环，故 J 包含 A 的极大理想的一个幂 (A, VIII, § 10, 第 1 节, 定理 1)。于是 \(\mathrm{prof}_A(\mathrm{J};\mathrm{C}_m)\geqslant \mathrm{prof}(\mathrm{C}_m) > m - p\) 对 \(m\geqslant p\) 成立（第 1 节, 命题 2 的推论 1），从而可以应用推论 1。

